\section{Special Functions}
\label{sec:specialfns}

The elementary functions---polynomials, the exponential, the trigonometric and
their inverses---run out long before mathematics does. The moment you write down
an integral like \(\int_0^\infty t^{z-1}e^{-t}\,dt\), or a power series whose
coefficients grow by a fixed ratio, or the solution of a second-order linear
differential equation with a variable coefficient, you have left the elementary
world and entered the realm of the \emph{special functions}: the gamma and zeta
functions, the error function, the Bessel and Airy functions, the Legendre and
hypergeometric families.\index{special functions} They are ``special'' only in
the sense that they were named and tabulated before computers existed; today they
are as ordinary as \(\sin\), and Mathilda treats them so.

What makes a special function useful in a computer algebra system is that it wears
three faces at once, and Mathilda shows you all three. There are the
\emph{exact symbolic values}---\(\Gamma(1/2)=\sqrt\pi\), \(\zeta(2)=\pi^2/6\)---which a
numeric library can only approximate. There are the \emph{symbolic identities}: the
recurrences, reflection and derivative rules, series expansions, and orthogonality
relations that let these functions participate in algebra rather than merely
returning numbers. And there is \emph{rigorous arbitrary-precision numerics}, for
when a number really is what you want, computed to any precision you ask for and,
where the kernel is backed by FLINT's ball arithmetic, with a guaranteed error
bound. This section takes the major families in turn, and the same three faces
recur in each: what closes into a clean symbol, what identity the system knows, and
how it puts a number on the value.

\subsection{The gamma family}
\label{ssec:sf-gamma}

Everything begins with the \B{Gamma} function, \(\Gamma(z)=\int_0^\infty
t^{z-1}e^{-t}\,dt\), the continuous interpolation of the
factorial.\index{gamma function} Its defining property, \(\Gamma(z+1)=z\,\Gamma(z)\)
with \(\Gamma(1)=1\), forces \(\Gamma(n)=(n-1)!\) at the integers, and from the
single exact value \(\Gamma(1/2)=\sqrt\pi\) the recurrence generates the whole
ladder of half-integers.

\begin{codepairs}
\pair{special-functions/gamma/1}{At the positive integers \B{Gamma} is the shifted
factorial: \(\Gamma(n)=(n-1)!\), so this is \(0!\) through \(5!\).}
\pair{special-functions/gamma/2}{The one transcendental seed value,
\(\Gamma(1/2)=\sqrt\pi\), returned exactly.}
\pair{special-functions/gamma/3}{The recurrence climbs from it:
\(\Gamma(7/2)=\tfrac{15}{8}\sqrt\pi\), a rational multiple of \(\sqrt\pi\).}
\pair{special-functions/gamma/4}{And descends to the negative half-integers---here
\(\Gamma(-1/2)=-2\sqrt\pi\), finite because the poles fall only on the
\emph{integers}.}
\pair{special-functions/gamma/5}{At the non-positive integers \(\Gamma\) has simple
poles, reported as \mcode{ComplexInfinity}.}
\pair{special-functions/gamma/6}{The two-argument form is the upper incomplete
gamma \(\Gamma(a,z)\); at a positive integer \(a\) it closes into elementary
form.}
\pair{special-functions/gamma/7}{With a numeric second argument that closed form is
an exact number: \(\Gamma(5,2)=168/e^{2}\).}
\pair{special-functions/gamma/8}{Where no symbol exists, \B{N} delivers digits---here
\(\Gamma(i)\) to thirty places, real and imaginary parts together.}
\pair{special-functions/gamma/9}{And \(\Gamma\) knows its own derivative, which
closes through the digamma function \(\psi=\Gamma'/\Gamma\).}
\end{codepairs}

The half-integer ladder (pairs~2--4) is the gamma function's signature, and it is
worth seeing why it terminates in \(\sqrt\pi\) rather than another special value.
\calloutnote{Theory}{The reflection formula
\(\Gamma(z)\,\Gamma(1-z)=\pi/\sin(\pi z)\)\index{gamma function!reflection formula}
evaluated at \(z=1/2\) gives \(\Gamma(1/2)^2=\pi\), hence \(\Gamma(1/2)=\sqrt\pi\).
Combined with the recurrence \(\Gamma(z+1)=z\,\Gamma(z)\), it pins every
half-integer to a rational multiple of \(\sqrt\pi\). The same reflection formula
explains the pole structure: since \(1/\Gamma\) is entire with zeros exactly at
\(0,-1,-2,\dots\), \(\Gamma\) itself has simple poles there and nowhere else, which
is why pair~4 is finite but pair~5 is not.} The numeric face (pairs~7--8) is the
one a differential-equation solver or a statistics routine actually leans
on.\calloutnote{Under the hood}{\B{Gamma} lives in
\texttt{src/special\_functions/gamma.c}. A machine-precision real routes to the C
library's \texttt{tgamma}; an arbitrary-precision real to MPFR's correctly rounded
\texttt{mpfr\_gamma}; and an arbitrary-precision \emph{complex} argument to Spouge's
approximation, whose coefficients are recomputed at runtime to the precision you
request, with the reflection formula handling \(\operatorname{Re}(z)<\tfrac12\). The
incomplete form uses \texttt{mpfr\_gamma\_inc} on the real axis and a lower-incomplete
series or a Lentz continued fraction in the complex plane.}

\usagebox{Gamma}

The \B{Beta} function \(B(a,b)=\Gamma(a)\Gamma(b)/\Gamma(a+b)\) is the gamma
function's closest relative---the normalizing constant of the beta distribution and
the value of the integral \(\int_0^1 t^{a-1}(1-t)^{b-1}\,dt\).\index{beta function}
Because it is built entirely from gammas, its exact reductions follow from theirs.

\begin{codepairs}
\pair{special-functions/beta/1}{A positive integer argument collapses the gamma
ratio to a rational: \(B(2,3)=1/12\).}
\pair{special-functions/beta/2}{Even with a rational partner it stays exact,
\(B(3,1/3)=27/14\).}
\pair{special-functions/beta/3}{Half-integers fold to rational multiples of \(\pi\)
instead---\(B(1/2,1/2)=\pi\), the arcsine normalization.}
\pair{special-functions/beta/4}{And \(B(5/2,7/2)=\tfrac{3}{256}\pi\).}
\pair{special-functions/beta/5}{When no such collapse is available, the residual
gamma form survives rather than being forced to a number.}
\pair{special-functions/beta/6}{A fully numeric call evaluates like any other.}
\end{codepairs}

Pair~5 is a small lesson in the system's temperament: \(B(1/3,1/3)\) is left as
\(\Gamma(1/3)^2/\Gamma(2/3)\) because no further exact simplification is available
without the reflection formula, which Mathilda does not apply automatically to a
product of gammas at complementary rationals.\calloutnote{Pitfall}{Do not expect
\(\Gamma(1/3)\,\Gamma(2/3)\) to collapse to \(2\pi/\sqrt3\) on its own. The gamma
reflection and multiplication identities are known to the theory but are \emph{not}
auto-applied---applying them blindly would fight as often as it helps---so products
of gammas at rational points are carried symbolically until you rewrite them by
hand. This is the honest cost of not over-simplifying.}

Two functions read the shape of the gamma function near its poles and along the
real axis. \B{LogGamma} is the log-gamma function \(\log\Gamma(z)\): the analytic
continuation of \(\log(\Gamma(z))\), not the same thing as \(\log(\Gamma(z))\)
term by term, and the right primitive whenever \(\Gamma\) itself would overflow.
\B{PolyGamma} is the family of logarithmic derivatives: \(\psi(z)=\psi^{(0)}(z)=
\Gamma'(z)/\Gamma(z)\) is the digamma function,\index{digamma function} and
\(\psi^{(n)}\) its \(n\)-th derivative.

\begin{codepairs}
\pair{special-functions/polygamma/1}{\B{LogGamma} keeps integers exact as the log of
a factorial: \(\log\Gamma(10)=\log(362880)=\log(9!)\).}
\pair{special-functions/polygamma/2}{At \(z=171\) the gamma function has already
overflowed a machine double---}
\pair{special-functions/polygamma/3}{---but \B{LogGamma} sails on, returning
\(\approx706.6\). This is why factorial ratios and log-likelihoods are computed
through it, never through \(\Gamma\) directly.}
\pair{special-functions/polygamma/4}{The digamma at a positive integer is a rational
minus Euler's constant: \(\psi(5)=\tfrac{25}{12}-\gamma\).}
\pair{special-functions/polygamma/5}{Gauss's digamma theorem closes rational
arguments into elementary form: \(\psi(1/2)=-\gamma-\log 4\).}
\pair{special-functions/polygamma/6}{The trigamma at \(1\) is \(\zeta(2)=\pi^2/6\).}
\pair{special-functions/polygamma/7}{\B{LogGamma} differentiates to digamma\ldots}
\pair{special-functions/polygamma/8}{\ldots and each further derivative raises the
polygamma order by one.}
\end{codepairs}

The contrast between pairs~2 and~3 is the entire reason \B{LogGamma} exists as a
first-class head rather than a shorthand for \texttt{Log[Gamma[z]]}: it stays finite
across a range where the gamma function runs from \(10^{-300}\) to \(10^{300}\), and
its imaginary part tracks the \emph{continued} argument past \(\pi\), where a naive
\(\log(\Gamma(z))\) would wrap.\calloutnote{Under the hood}{\B{PolyGamma}
(\texttt{src/special\_functions/polygamma.c}) handles order \(n=0\) at rational
arguments through the Gauss digamma theorem---a finite trigonometric sum---and closes
odd orders through \(\zeta(n{+}1)\). Numerically it uses MPFR's digamma for
\(n=0\) and a recurrence-shift plus a Bernoulli asymptotic expansion for higher
orders; complex arbitrary-precision values are handed to FLINT's \texttt{acb\_polygamma}.
The negative order \(\psi^{(-1)}(z)\) is \B{LogGamma} itself.}

The last member of the family is the \B{Pochhammer} symbol \((a)_n=
a(a+1)\cdots(a+n-1)=\Gamma(a+n)/\Gamma(a)\), the \emph{rising factorial} that
supplies the coefficients of every hypergeometric series we meet at the end of this
section.\index{Pochhammer symbol}

\begin{codepairs}
\pair{special-functions/pochhammer/1}{A symbolic base gives an explicit polynomial
product of \(n\) linear factors.}
\pair{special-functions/pochhammer/2}{A numeric base gives an exact value,
\((10)_6=3603600\).}
\pair{special-functions/pochhammer/3}{\((1)_{25}\) is just \(25!\), carried exactly
as a GMP big integer.}
\pair{special-functions/pochhammer/4}{Half-integer arguments reduce through the gamma
ratio: \((1/2)_3=15/8\).}
\pair{special-functions/pochhammer/5}{And \((3/2)_{1/2}=2/\sqrt\pi\), a gamma ratio
that happens to close.}
\end{codepairs}

\subsection{The zeta family}
\label{ssec:sf-zeta}

The Riemann zeta function \(\zeta(s)=\sum_{n\ge1}n^{-s}\) is the deepest object in
this section, and \B{Zeta} carries its exact arithmetic.\index{Riemann zeta
function} At the even positive integers it closes into powers of \(\pi\); its values
at the non-positive integers, defined by analytic continuation, are rational; and
its values at the odd positive integers are believed to be new transcendentals with
no closed form at all.

\begin{codepairs}
\pair{special-functions/zeta/1}{The even values are rational multiples of powers of
\(\pi\): \(\zeta(2)=\pi^2/6\)\ldots}
\pair{special-functions/zeta/2}{\ldots and \(\zeta(4)=\pi^4/90\).}
\pair{special-functions/zeta/3}{This is exactly the Basel sum, which \B{Sum}
recognizes independently and returns identically.}
\pair{special-functions/zeta/4}{Analytic continuation gives \(\zeta(0)=-1/2\)\ldots}
\pair{special-functions/zeta/5}{\ldots and the famous \(\zeta(-1)=-1/12\).}
\pair{special-functions/zeta/6}{The negative even integers are the \emph{trivial
zeros}: \(\zeta(-2)=0\).}
\pair{special-functions/zeta/7}{The odd values have no known closed form, so
\(\zeta(3)\) is returned as itself---a genuine constant, not a failure.}
\pair{special-functions/zeta/8}{But it is a \emph{number}: \B{N} gives Apéry's
constant to thirty digits.}
\pair{special-functions/zeta/9}{On the critical line, at the height \(t=14.1347\ldots\)
of the first nontrivial zero, \(\zeta(1/2+it)\) collapses to the order of
\(10^{-7}\)---the Riemann hypothesis made visible.}
\pair{special-functions/zeta/10}{The Hurwitz zeta \(\zeta(s,a)=\sum_{n\ge0}(n+a)^{-s}\)
shifts the denominators, and still closes: \(\zeta(2,3)=\pi^2/6-5/4\).}
\pair{special-functions/zeta/11}{The Stieltjes constants are inert symbols; the
zeroth is Euler's constant \(\gamma\).}
\end{codepairs}

Pairs~4--6 are not sleight of hand but the content of the \emph{functional
equation}, and pair~9 is a numerical glimpse of the most famous open problem in
mathematics.\calloutnote{Theory}{Riemann's functional
equation\index{Riemann zeta function!functional equation}
\(\zeta(s)=2^s\pi^{s-1}\sin(\pi s/2)\,\Gamma(1-s)\,\zeta(1-s)\) relates the value at
\(s\) to the value at \(1-s\), and with it the whole left half-plane is determined by
the right. The factor \(\sin(\pi s/2)\) vanishes at the negative even integers,
producing the trivial zeros of pair~6; the values \(\zeta(0)=-1/2\) and
\(\zeta(-1)=-1/12\) fall out of the same equation. The \emph{nontrivial} zeros all
lie in the strip \(0<\operatorname{Re}(s)<1\), and the \emph{Riemann hypothesis}
\index{Riemann hypothesis} asserts they lie exactly on the line
\(\operatorname{Re}(s)=1/2\); pair~9 evaluates \(\zeta\) at the first of them, where
it is zero to the precision requested.} The even values, by contrast, are elementary
once you know the Bernoulli numbers.\calloutnote{Theory}{Euler's formula
\(\zeta(2n)=(-1)^{n+1}\dfrac{B_{2n}(2\pi)^{2n}}{2\,(2n)!}\) expresses every even zeta
value through a \B{BernoulliB} number\index{Bernoulli numbers}; for \(n=1\),
\(B_2=1/6\) gives \(\zeta(2)=\pi^2/6\). No analogous formula is known for the odd
values, which is why \(\zeta(3)\) (pair~7) has no closed form while \(\zeta(2)\) and
\(\zeta(4)\) do.} Underneath the exact layer is a rigorous numeric
one.\calloutnote{Under the hood}{\B{Zeta}, \B{HurwitzZeta} and \B{StieltjesGamma}
(\texttt{src/special\_functions/zeta.c}, and \texttt{src/flint\_num\_bridge.c} for
the numerics) delegate arbitrary-precision \emph{complex} evaluation to FLINT's Arb
library---\texttt{acb\_dirichlet\_zeta}, \texttt{acb\_dirichlet\_hurwitz},
\texttt{acb\_dirichlet\_stieltjes}. Arb works in \emph{ball arithmetic}\index{ball
arithmetic}: every value carries a midpoint and a rigorous radius, so pair~9's
near-zero comes with a proof that the true value lies within a known ball of the
printed digits---not a hope that rounding behaved.}

\usagebox{Zeta}

The Stieltjes constants \(\gamma_n\) (pair~11) are the coefficients of the Laurent
expansion of \(\zeta\) about its pole at \(s=1\); the zeroth is Euler's constant
\(\gamma\), surfaced here as \B{EulerGamma}, and the rest are carried as inert
\B{StieltjesGamma} symbols.

\subsection{Error functions and the integral functions}
\label{ssec:sf-erf}

The error function \(\operatorname{erf}(z)=\frac{2}{\sqrt\pi}\int_0^z e^{-t^2}\,dt\)
is the antiderivative of the Gaussian, and no elementary function equals
it---which is precisely why it needs a name.\index{error function} \B{Erf}, its
complement \B{Erfc}, the imaginary variant \B{Erfi}, and the inverse \B{InverseErf}
form a small closed world.

\begin{codepairs}
\pair{special-functions/erf/1}{\(\operatorname{erf}(0)=0\)\ldots}
\pair{special-functions/erf/2}{\ldots and \(\operatorname{erf}(\infty)=1\): the total
mass of the normalized Gaussian.}
\pair{special-functions/erf/3}{It is an odd function, and the odd symmetry is applied
symbolically.}
\pair{special-functions/erf/4}{Arbitrary-precision numerics track the digits you ask
for.}
\pair{special-functions/erf/5}{The complement \(\operatorname{erfc}=1-\operatorname{erf}\)
vanishes at infinity, computed cancellation-free so the tail keeps its
significance.}
\pair{special-functions/erf/6}{The imaginary error function
\(\operatorname{erfi}(z)=-i\,\operatorname{erf}(iz)\) has its own kernel.}
\pair{special-functions/erf/7}{The derivative is the Gaussian itself---so the whole
Taylor series is generated from this one rule.}
\pair{special-functions/erf/8}{The inverse function has exact endpoints:
\(\operatorname{erf}^{-1}(1)=\infty\).}
\pair{special-functions/erf/9}{And numeric values in between, on \([-1,1]\).}
\end{codepairs}

The split between \B{Erf} and \B{Erfc} (pairs~2 and~5) is not redundancy but
numerics: for large positive \(z\), \(\operatorname{erf}(z)\) is so close to \(1\)
that \(1-\operatorname{erf}(z)\) loses every significant digit to cancellation, so
\(\operatorname{erfc}\) is computed by its own cancellation-free
route.\calloutnote{Under the hood}{The kernels live in
\texttt{src/special\_functions/erf.c} (and \texttt{inverf.c} for the inverse). Real
arguments use the C library or MPFR's \texttt{erf}/\texttt{erfc}; complex arguments
sum the cancellation-aware Maclaurin series of DLMF~7.6.2 in MPFR with
\(|z|^2/\ln2\) guard bits, so even a machine-precision complex result carries full
accuracy. Neither the C library nor MPFR ships an inverse error function, so
\B{InverseErf} is Newton's iteration on \(\operatorname{erf}(z)-s\), started from a
Winitzki seed and polished to the working precision.}

The same principle---name the antiderivative when it is not elementary---generates a
whole class of \emph{integral functions}, and Mathilda produces them directly when
you integrate.

\begin{codepairs}
\pair{special-functions/integrals/1}{The sine integral \(\operatorname{Si}\) is the
antiderivative of the sinc function, and \B{Integrate} returns it by name.}
\pair{special-functions/integrals/2}{The cosine integral \(\operatorname{Ci}\)
likewise.}
\pair{special-functions/integrals/3}{The exponential integral \(\operatorname{Ei}\)
names \(\int e^x/x\,dx\).}
\pair{special-functions/integrals/4}{And \(\int dx/\log x\) is the same function in
disguise---the logarithmic integral \(\operatorname{li}(x)=\operatorname{Ei}(\log x)\).}
\end{codepairs}

Each of these is a genuine function with the full three-faced treatment: it
evaluates, it differentiates, it has special values.

\begin{codepairs}
\pair{special-functions/integrals/5}{\B{SinIntegral} to twenty digits\ldots}
\pair{special-functions/integrals/6}{\ldots \B{CosIntegral}\ldots}
\pair{special-functions/integrals/7}{\ldots \B{ExpIntegralEi}\ldots}
\pair{special-functions/integrals/8}{\ldots and \B{LogIntegral}, whose growth rate is
the density of the primes.}
\pair{special-functions/integrals/9}{\(\operatorname{Si}\) saturates to \(\pi/2\) at
infinity---the Dirichlet integral \(\int_0^\infty \sin(x)/x\,dx\).}
\pair{special-functions/integrals/10}{The Fresnel integral \(C\) of optics and
diffraction, \B{FresnelC}\ldots}
\pair{special-functions/integrals/11}{\ldots and its partner \B{FresnelS}.}
\pair{special-functions/integrals/12}{And \(\operatorname{Si}\) differentiates back
to the cardinal sine \B{Sinc}---the loop closes.}
\end{codepairs}

That \(\int dx/\log x\) and \(\int e^x/x\,dx\) come back as the \emph{same} function
(pairs~3--4) is the logarithmic integral's claim to fame: \(\operatorname{li}(x)\)
is the leading term of the prime-counting function, and the substitution
\(t=\log x\) turns one integral into the other.\calloutnote{Under the hood}{On the
real axis \B{ExpIntegralEi} uses MPFR's correctly rounded \texttt{mpfr\_eint} for
\(x>0\) and an on-cut convergent series for \(x<0\); large \(|z|\) switches to the
asymptotic expansion, and the two regimes overlap and agree. \B{LogIntegral} is
literally \(\operatorname{Ei}(\log z)\), reusing the same kernel. The sine, cosine,
and Fresnel integrals sum their Maclaurin series for moderate arguments and their
trigonometric-prefactored asymptotic forms for large ones, with the complex plane
handled by the shared MPFR-complex toolkit.}

\subsection{Bessel and Airy functions}
\label{ssec:sf-bessel}

Bessel functions are what solutions to the wave and heat equations look like in
cylindrical coordinates, and they are the archetype of a function defined by a
differential equation rather than a formula.\index{Bessel functions} \B{BesselJ}
is the function of the first kind, regular at the origin; \B{BesselY} the second
kind, singular there; and \B{BesselI}, \B{BesselK} the modified pair for the
hyperbolic version of the equation.

\begin{codepairs}
\pair{special-functions/bessel/1}{\(J_0(0)=1\); every higher \(J_n\) vanishes at the
origin.}
\pair{special-functions/bessel/2}{At half-integer order the Bessel function is
\emph{elementary}: \(J_{1/2}(x)=\sqrt{2/\pi x}\,\sin x\).}
\pair{special-functions/bessel/3}{And \(J_{-1/2}(x)=\sqrt{2/\pi x}\,\cos x\)---the
spherical Bessel functions in disguise.}
\pair{special-functions/bessel/4}{Numerics to any precision.}
\pair{special-functions/bessel/5}{The Frobenius series about the origin, generated on
demand.}
\pair{special-functions/bessel/6}{The recurrence shows itself in the derivative:
\(J_0'=-J_1\).}
\pair{special-functions/bessel/7}{The second kind \(Y_0\), singular at the
origin\ldots}
\pair{special-functions/bessel/8}{\ldots the modified \(I_0\)\ldots}
\pair{special-functions/bessel/9}{\ldots and \(K_0\), the exponentially decaying
modified solution.}
\end{codepairs}

The half-integer reductions (pairs~2--3) are a real theorem, not a lookup: they are
why a quantum-mechanics problem in a spherical well produces \(\sin\) and \(\cos\)
over \(\sqrt r\).\calloutnote{Theory}{The Bessel functions \(J_\nu\) and \(Y_\nu\)
are the two independent solutions of Bessel's equation
\(x^2y''+xy'+(x^2-\nu^2)y=0\). When \(\nu\) is a half-integer the equation degenerates
just enough for the solution to close in elementary terms: \(J_{1/2}\) and
\(J_{-1/2}\) are \(\sin\) and \(\cos\) scaled by \(\sqrt{2/\pi x}\), and the
three-term recurrence \(J_{\nu-1}(x)+J_{\nu+1}(x)=\tfrac{2\nu}{x}J_\nu(x)\) generates
every other half-integer order from them. Mathilda applies exactly this reduction
(\texttt{src/special\_functions/bessel.c}).}

\usagebox{BesselJ}

The Airy functions are the simplest functions that oscillate on one side of the
origin and decay on the other, the model for a wave near a turning
point.\index{Airy functions} \B{AiryAi} and \B{AiryBi} solve \(y''=xy\), and their
values at the origin tie the whole family back to the gamma function.

\begin{codepairs}
\pair{special-functions/airy/1}{\(\operatorname{Ai}(0)\) closes through the gamma
function: \(1/(3^{2/3}\Gamma(2/3))\).}
\pair{special-functions/airy/2}{\(\operatorname{Bi}(0)\) shares the same
\(\Gamma(2/3)\), scaled differently.}
\pair{special-functions/airy/3}{The derivative at the origin brings in
\(\Gamma(1/3)\) instead.}
\pair{special-functions/airy/4}{The Airy value numerically, to thirty digits\ldots}
\pair{special-functions/airy/5}{\ldots and at \(x=1\), already deep in the decaying
region.}
\pair{special-functions/airy/6}{\B{AiryAi} differentiates to its own derivative head,
\B{AiryAiPrime}, closing the family under \B{D}.}
\end{codepairs}

That \(\operatorname{Ai}(0)\) and \(\operatorname{Ai}'(0)\) involve \(\Gamma(2/3)\)
and \(\Gamma(1/3)\) respectively (pairs~1--3) is the thread tying this subsection
back to the first: the special functions are a single web, and the gamma function
sits near its center.\calloutnote{Theory}{The Airy equation \(y''=xy\) has no
singular points at finite \(x\), so both solutions are entire. The initial values
\(\operatorname{Ai}(0)=3^{-2/3}/\Gamma(2/3)\) and
\(\operatorname{Ai}'(0)=-3^{-1/3}/\Gamma(1/3)\) come from the Mellin-transform
representation of the solution, and every Taylor coefficient follows from the
equation by \(y^{(n+2)}=x\,y^{(n)}+n\,y^{(n-1)}\). The kernels are in
\texttt{src/special\_functions/airyai.c} and \texttt{airybi.c}.}

\subsection{Legendre functions of the first and second kind}
\label{ssec:sf-legendre}

Legendre's equation \((1-x^2)y''-2xy'+n(n+1)y=0\) is the angular part of Laplace's
equation on the sphere, and it too has two independent
solutions.\index{Legendre functions} \B{LegendreP} is the first: at a non-negative
integer degree it is the Legendre \emph{polynomial}, the workhorse of Gaussian
quadrature and multipole expansions. \B{LegendreQ} is the second kind---the other
solution of the same equation---and it is new in this release.\index{Legendre
functions!second kind}

\begin{codepairs}
\pair{special-functions/legendre/1}{The Legendre polynomials \(P_0\) through \(P_3\),
generated on demand.}
\pair{special-functions/legendre/2}{\(P_4\), with its characteristic rational
coefficients.}
\pair{special-functions/legendre/3}{Orthogonality on \([-1,1]\): the squared norm is
\(\int_{-1}^1 P_n^2 = 2/(2n{+}1)\), here \(2/5\).}
\pair{special-functions/legendre/4}{And distinct degrees are orthogonal---the
integral of \(P_2P_3\) is exactly zero.}
\pair{special-functions/legendre/5}{The second-kind solution \(Q_0\) is the
logarithm \(\tfrac12\log\frac{1+x}{1-x}=\operatorname{arctanh}x\).}
\pair{special-functions/legendre/6}{\(Q_1\) carries that same logarithm, times
\(x\), minus one.}
\pair{special-functions/legendre/7}{\(Q_2\) pairs the logarithm with the polynomial
\(P_2\)---the general pattern.}
\pair{special-functions/legendre/8}{And \(Q\) evaluates numerically like any other
function.}
\end{codepairs}

Pairs~3--4 are the property that makes the Legendre polynomials indispensable, and
the system computes them exactly rather than numerically.\calloutnote{Theory}{The
polynomials \(P_n\) are orthogonal on \([-1,1]\):
\(\int_{-1}^1 P_m(x)P_n(x)\,dx = \frac{2}{2n+1}\delta_{mn}\). This single relation is
what lets any function on the interval be expanded in a Legendre series, and it is the
foundation of Gauss--Legendre quadrature. The second-kind solution \(Q_n\) satisfies
the same differential equation but is singular at \(x=\pm1\)---the endpoints where the
logarithm in pairs~5--7 blows up---so a physical solution regular on the whole sphere
keeps only the \(P_n\) part. \B{LegendreQ} follows the pattern
\(Q_n(x)=P_n(x)\,Q_0(x)-W_{n-1}(x)\) with \(Q_0=\operatorname{arctanh}x\); the
implementation is in \texttt{src/special\_functions/legendre.c}.} The visible
structure of \(Q_2\) in pair~7---a Legendre polynomial multiplying
\(\operatorname{arctanh}x\), minus a lower-degree polynomial---is exactly that
recurrence made concrete.

\subsection{The hypergeometric family: one function to rule them all}
\label{ssec:sf-hypergeometric}

Almost every function in this section is a special case of one
master series.\index{hypergeometric function} The generalized hypergeometric
function
\[
  {}_pF_q\!\left(\begin{matrix}a_1,\dots,a_p\\ b_1,\dots,b_q\end{matrix};z\right)
  = \sum_{k=0}^\infty
    \frac{(a_1)_k\cdots(a_p)_k}{(b_1)_k\cdots(b_q)_k}\,\frac{z^k}{k!}
\]
is a power series whose term ratio is a fixed rational function of the summation
index---and that one structural constraint is general enough to contain the
exponential, the logarithm, the trigonometric functions, the Bessel functions, the
error function, and the Legendre functions as special cases. Mathilda exposes it as
\B{HypergeometricPFQ}, with the classical low-order cases \B{Hypergeometric2F1}
(Gauss), \B{Hypergeometric1F1} (Kummer's confluent), and \B{Hypergeometric0F1}
(the confluent limit) as named heads. The system's job is to \emph{recognize} when a
hypergeometric collapses to something elementary, and it does.

\begin{codepairs}
\pair{special-functions/hypergeometric/1}{Gauss's \({}_2F_1(1,1;2;x)\) is a
logarithm in disguise: \(-\log(1-x)/x\).}
\pair{special-functions/hypergeometric/2}{A negative-integer parameter terminates the
series into a polynomial---here \((1-x)^3\).}
\pair{special-functions/hypergeometric/3}{Kummer's \({}_1F_1(1;2;x)\) is
\((e^x-1)/x\)\ldots}
\pair{special-functions/hypergeometric/4}{\ldots and when its two parameters agree it
is the exponential itself.}
\pair{special-functions/hypergeometric/5}{The confluent \({}_0F_1(;1/2;x)\) is a
hyperbolic cosine of a square root.}
\pair{special-functions/hypergeometric/6}{The generalized \({}_pF_q\) carries any
number of parameters, and reduces by the same rules---here matching pair~1.}
\pair{special-functions/hypergeometric/7}{Where no elementary form exists, it is a
number: this one is \(2\log 2\).}
\end{codepairs}

The unification is not just aesthetic; it is how a CAS avoids re-implementing the
same reductions a dozen times.\calloutnote{Theory}{A function is
\emph{hypergeometric} when the ratio of consecutive series coefficients,
\(c_{k+1}/c_k\), is a rational function of \(k\)---and astonishingly many functions
are. The exponential is \({}_0F_0\); \(\log(1-x)=-x\,{}_2F_1(1,1;2;x)\) (pair~1);
\(J_\nu(x)\) is a \({}_0F_1\) up to a power factor (compare pair~5's \(\cosh\), which
is \(J_{-1/2}\) territory); \(\operatorname{erf}\), the Legendre and Gegenbauer
polynomials, and the Airy functions are all \({}_pF_q\) with particular parameters.
The two levers that produce a closed form are visible above: a \emph{negative
integer} upper parameter truncates the infinite series to a polynomial (pair~2), and
special parameter relations trigger classical summation theorems (pairs~1, 3--5).
Mathilda's reductions live in \texttt{src/special\_functions/hypergeopfq.c}.}

Three functions extend the reach of the elementary world one step further and round
out the family. The polylogarithm \B{PolyLog} generalizes the logarithm to a whole
ladder of orders;\index{polylogarithm} the Lerch transcendent \B{LerchPhi} contains
the polylogarithm, the Hurwitz zeta, and much else as special
cases;\index{Lerch transcendent} and the Lambert \(W\) function \B{ProductLog}
inverts \(w\,e^w=z\), the transcendental equation that appears whenever an unknown
sits both inside and outside an exponential.\index{Lambert W function}

\begin{codepairs}
\pair{special-functions/polylog/1}{Order one is the logarithm:
\(\operatorname{Li}_1(x)=-\log(1-x)\).}
\pair{special-functions/polylog/2}{The dilogarithm at \(1\) is \(\zeta(2)=\pi^2/6\).}
\pair{special-functions/polylog/3}{At \(-1\) it is \(-\pi^2/12\).}
\pair{special-functions/polylog/4}{And the celebrated closed form
\(\operatorname{Li}_2(1/2)=\pi^2/12-\tfrac12\log^2 2\).}
\pair{special-functions/polylog/5}{At argument \(1\) the order-\(s\) polylog is
\(\zeta(s)\), so \(\operatorname{Li}_3(1)=\zeta(3)\).}
\pair{special-functions/polylog/6}{The Lerch transcendent reduces to the same
logarithm for these arguments---the master function containing all the rest.}
\pair{special-functions/polylog/7}{The Lambert \(W\): \(W(e)=1\), since
\(1\cdot e^1=e\).}
\pair{special-functions/polylog/8}{And \(W(-1/e)=-1\), the branch point.}
\pair{special-functions/polylog/9}{Its value at \(1\) is the omega constant,
\(0.5671\ldots\).}
\pair{special-functions/polylog/10}{Its derivative is written implicitly in terms of
itself---the signature of a function defined by inversion.}
\end{codepairs}

The dilogarithm identities (pairs~2--4) are among the most beautiful closed forms in
analysis, and that Mathilda knows \(\operatorname{Li}_2(1/2)=\pi^2/12-\tfrac12\log^2
2\) exactly rather than as a decimal is the whole point of a symbolic
system.\calloutnote{Under the hood}{\B{PolyLog}
(\texttt{src/special\_functions/polylog.c}), \B{LerchPhi} (\texttt{lerchphi.c}), and
\B{ProductLog} (\texttt{productlog.c}) each carry their own special-value tables and
their own numeric kernels. The Lerch transcendent
\(\Phi(z,s,a)=\sum_{k\ge0}z^k/(k+a)^s\) is the genuine master function of the three:
\(\operatorname{Li}_s(z)=z\,\Phi(z,s,1)\) and \(\zeta(s,a)=\Phi(1,s,a)\), so pair~6's
reduction and the Hurwitz zeta of \S\ref{ssec:sf-zeta} are two faces of one series.
\B{ProductLog}'s implicit derivative (pair~10) comes from differentiating
\(w\,e^w=z\) directly, since no explicit formula for \(W\) exists.}

\subsection*{Where this connects}

The special functions are the point where the three layers of Mathilda meet on a
single object. Their \emph{exact} values are computed by the same rational and
symbolic arithmetic as \S\ref{sec:arithmetic} and \S\ref{sec:algebra}; their
\emph{identities}---the derivatives that close the gamma, Bessel, Airy, and Legendre
families under \B{D}, the series expansions, the antiderivatives that \B{Integrate}
returns by name---are what let them take part in the calculus of
\S\ref{sec:calculus} rather than merely being evaluated at its end; and their
\emph{numerics}, backed by MPFR on the real line and by FLINT's rigorous Arb ball
arithmetic in the complex plane, are what the numerical methods of
\S\ref{sec:numcalc} call when an integral or a differential equation resolves into
one of these functions. The orthogonality of the Legendre polynomials underwrites
Gaussian quadrature; the gamma and zeta functions thread through analytic number
theory; and the hypergeometric family is the organizing idea that keeps the whole
collection from being a bag of unrelated formulas. For the exhaustive options and
corner cases of any function named here---the branch-cut conventions, the incomplete
and generalized forms, the full list of recognized reductions---its reference page
is one \texttt{?Name} away at the prompt.
